% Generated by scripts/assemble.py; corrections belong in source/corrections.json.
\examhead{163}{广西大学 2025 年数学分析试卷}{155}
% p156-b10; source PDF page 156
\par\Needspace{5\baselineskip}
\blockanchor{p156-b10}
\majorquestion{一}填空题. 每题 5 分, 共 25 分.\par
% p156-b11; source PDF page 156
\blockanchor{p156-b11}
\question{1}$\lim_{x\to0}\frac{1}{x}\int_{0}^{x}\cos t^{3}dt=$  \_\_\_\_.\par
% p156-b12; source PDF page 156
\blockanchor{p156-b12}
\question{2}$f(x)=\frac{3x^{3}+4}{x^{2}-2x}$  的渐近线为 \_\_\_\_.\par
% p156-b13; source PDF page 156
\blockanchor{p156-b13}
\question{3}幂级数  $\sum_{n=1}^{\infty}\frac{(x-2)^{2n-1}}{(2n-1)!}$  的收敛区域为 \_\_\_\_.\par
% p156-b14; source PDF page 156
\blockanchor{p156-b14}
\question{4}u = xyz 在点  $(5,1,2)$  沿点  $(5,1,2)$  到  $(9,4,14)$  的方向导数为 \_\_\_\_.\par
% p156-b15; source PDF page 156
\blockanchor{p156-b15}
\question{5}L 是球面  $x^{2} + y^{2} + z^{2} = a^{2}$  与 x = y 相交的圆周，则  $\int_{L} \sqrt{2y^{2} + z^{2}} \, ds =$ \_\_\_\_ .\par
% p156-b16; source PDF page 156
\par\Needspace{5\baselineskip}
\blockanchor{p156-b16}
\majorquestion{二}计算题. 每题 10 分, 共 60 分.\par
% p156-b17; source PDF page 156
\blockanchor{p156-b17}
\question{6}求 $y = x^{2}$ 与 $y = 8 - x^{2}$ 所围区域的面积.\par
% p156-b18; source PDF page 156
\blockanchor{p156-b18}
\question{7}$u = x^{2} + y^{2} + z^{2}$, 其中 $z = f(x, y)$ 是由 $x^{3} + y^{3} + z^{3} = 3xyz$ 确定的隐函数, 求 $u_{x}, u_{xx}$.\par
% p156-b19; source PDF page 156
\par\Needspace{6\baselineskip}
\blockanchor{p156-b19}
\question{8}设 $f(x) = \begin{cases} -\frac{\pi}{4}, & -\pi < x < 0; \\ \frac{\pi}{4}, & 0 \leq x < \pi. \end{cases}$ 求 $f(x)$ 的傅里叶级数展开式. 并推出\par
\nopagebreak[4]
% p156-b20; source PDF page 156
\blockanchor{p156-b20}
\displayformula{\frac {\sqrt {3}}{6} \pi = 1 - \frac {1}{5} + \frac {1}{7} - \frac {1}{11} + \frac {1}{13} - \frac {1}{17} + \dots .}
% p156-b21; source PDF page 156
\blockanchor{p156-b21}
\question{9}求 $(x_0, y_0, z_0)$ 到平面 $Ax + By + Cz + D = 0$ 最短距离.\par
% p156-b22; source PDF page 156
\blockanchor{p156-b22}
\question{10}求 $I = \iiint_{V} z^{2} \, \mathrm{d}x \, \mathrm{d}y \, \mathrm{d}z$, $V$ 是由锥面 $z^{2} = \frac{h^{2}}{R^{2}} (x^{2} + y^{2})$ 与平面 $z = h$ 围成的区域.\par
% p156-b23; source PDF page 156
\par\Needspace{6\baselineskip}
\blockanchor{p156-b23}
\question{11}求曲面积分\par
\nopagebreak[4]
% p156-b24; source PDF page 156
\blockanchor{p156-b24}
\displayformula{I = \iint_ {S} y (x - z) \mathrm{d} y \mathrm{d} z + x ^ {2} \mathrm{d} z \mathrm{d} x + (y ^ {2} + z x) \mathrm{d} x \mathrm{d} y.}
% p156-b25; source PDF page 156
\blockanchor{p156-b25}
其中 $S$ 是曲面 $z = 5 - x^{2} - y^{2}$ 上 $z \geq 1$ 部分的外侧.\par
% p156-b26; source PDF page 156
\par\Needspace{5\baselineskip}
\blockanchor{p156-b26}
\majorquestion{三}证明题. 每题 10 分, 共 40 分.\par
% p157-b00; source PDF page 157
\blockanchor{p157-b00}
\question{12}$\lim_{n\to \infty}a_n = a,a > 0,$ 证明： $\lim_{n\to \infty}\sqrt[n]{a_n} = 1.$\par
% p157-b01; source PDF page 157
\blockanchor{p157-b01}
\question{13}证明 $f(x) = x^2$ 在 $[a, b]$ 上一致连续，在 $(- \infty, +\infty)$ 上不一致连续.\par
% p157-b02; source PDF page 157
\par\Needspace{6\baselineskip}
\blockanchor{p157-b02}
\question{14}设曲线\par
\nopagebreak[4]
% p157-b03; source PDF page 157
\blockanchor{p157-b03}
\displayformula{\left\{ \begin{array}{l} x = 2 (\sin t - t \cos t); \\ y = 2 (\cos t + t \sin t). \end{array} \right.}
% p157-b04; source PDF page 157
\blockanchor{p157-b04}
证明原点到曲线上任意一点处的法线的距离为常数.\par
% p157-b05; source PDF page 157
\blockanchor{p157-b05}
\question{15}根据有限覆盖定理证明: 实数轴上的任一有界无限点集至少有一个聚点.
\majorquestion{四}综合题. 分值未知.\par
% p157-b06; source PDF page 157
\par\Needspace{6\baselineskip}
\blockanchor{p157-b06}
\question{16}用\par
\nopagebreak[4]
% p157-b07; source PDF page 157
\blockanchor{p157-b07}
\displayformula{f (x, y) = \left\{ \begin{array}{l l} (x ^ {2} + y ^ {2}) \sin \frac {1}{\sqrt {x ^ {2} + y ^ {2}}}, & x ^ {2} + y ^ {2} \neq 0; \\ 0, & x ^ {2} + y ^ {2} = 0. \end{array} \right.}
% p157-b08; source PDF page 157
\blockanchor{p157-b08}
说明偏导数连续并不是函数可微的必要条件.\par
% p157-b09; source PDF page 157
\blockanchor{p157-b09}
\question{17}用求导, 积分计算 $I(x) = \int_{0}^{+\infty} e^{-t^2} \cos 2xt \, \mathrm{d}t$, 并说明原因.\par
